\section{Des cas aux archétypes~: ce que l'histoire révèle et ce qu'elle ne peut pas expliquer}\label{sec:reduction}
% =============================================================================

Le parcours qui précède a traversé dix cas sur deux siècles, distribués entre les sections §\,\ref{sec:era0} (avant l'électricité) et §\,\ref{sec:contemporain} (systèmes informationnels contemporains). Du télégraphe optique de Chappe au datacenter d'IA générative, du calcul à la plume de Clairaut et Lalande à l'inférence par GPU à 1\,400~watts, de la gutta-percha du câble transatlantique au bouquet élémentaire des fonderies de semi-conducteurs, le matériau est hétérogène par nature et par époque. La méthode a été inductive~: chaque cas a été observé pour lui-même, les catégories analytiques ont émergé de la comparaison entre les cas, et les patterns provisoires qu'un cas faisait apparaître ont été modifiés ou infirmés par le suivant. Le moment est venu de faire le geste inverse~: non plus observer des cas, mais rassembler ce que les observations révèlent de commun, et nommer ce qu'elles ne parviennent pas à expliquer.

La répartition des cas dans la synthèse n'est pas uniforme. Les régularités qui apparaissent dans la première partie de ce bilan s'établissent sur le matériau du XIX\textsuperscript{e} et du début du XX\textsuperscript{e} siècle, où plusieurs cas comparables permettent l'induction. Les cas contemporains (§\,\ref{sec:computing}, §\,\ref{sec:convergence}, §\,\ref{sec:contemporain}) interviennent plus tard dans le bilan~: ce sont eux qui font apparaître l'\emph{anomalie} énergétique et les deux postulats qui cessent de tenir. Les premiers servent à identifier ce qui se répète~; les seconds, à identifier ce qui rompt. La double fonction est délibérée.

\subsection{Trois régularités et une anomalie}\label{subsec:regularites}

La première régularité est institutionnelle. Chaque réseau de communication observé dans ce chapitre converge vers une forme ou une autre de monopole~: sept trajectoires nationales pour le télégraphe (§\,\ref{sec:telegraph}), quatre pour le téléphone (§\,\ref{sec:telephone}), trois mécanismes successifs pour le wireless (§\,\ref{sec:wireless}), le monopole régulé d'Édison et d'Insull pour l'électrification (§\,\ref{sec:electrification}), la domination d'IBM à 85,7\,\% du parc mécanographique (§\,\ref{sec:mechanography}). Les conditions initiales varient radicalement~--- sécurité ferroviaire, souveraineté militaire, contrôle policier, modernisation autoritaire, concession coloniale~--- mais le résultat institutionnel est le même. Ce qui varie, c'est le mécanisme~: rendements de réseau pour le télégraphe, verrouillage du consommable pour la tabulatrice, intégration verticale et contrôle de l'interconnexion pour le téléphone, brevets et refus d'interopérabilité pour Marconi, rendements d'échelle dans la génération pour l'électricité. La forme apparente est récurrente~; les mécanismes qui la produisent ne le sont pas. La concentration des infrastructures informationnelles est donc une régularité morphologique sans mécanisme commun, et non un résultat causalement unifié~: ce constat oblige le Ch.~\ref{ch:structures} à la décomposer en dimensions séparées plutôt qu'à la traiter comme un mécanisme à part entière, dans les rendements croissants (D2), la nature du capital (D3), la coordination par les standards (D7) et l'allocation politique des ressources rares comme le spectre (D6).

L'envers de cette régularité institutionnelle est tout aussi instructif. Babbage occupe la place du cas qui ne devient pas infrastructure malgré la conjonction d'éléments qui auraient pu en faire un. Une demande de calcul vérifiable existe~: les tables du Nautical Almanac contiennent des erreurs identifiées, et Prony a démontré que la division du travail intellectuel pouvait produire des résultats reproductibles. Une capacité technique existe~: les plans de la Difference Engine sont défendables, ceux de l'Analytical Engine plus exigeants mais conceptuellement cohérents. Pourtant aucune des deux ne deviendra machine de production. Ce qui manque est identifiable~: une organisation capable d'absorber les coûts d'une fabrication mécanique de précision sans débouché commercial avéré, un véhicule financier adapté à un projet dont les bénéfices ne sont pas appropriables par le financeur, une promesse stabilisée auprès des utilisateurs administratifs qui auraient pu être les premiers clients. Babbage n'est pas un génie en avance sur son temps~; c'est un cas de désajustement entre projet technique, demande solvable, administration du financement et trajectoire d'usage\footnote{La formule \og en avance sur son temps\fg{} déplace le problème vers une métaphysique chronologique sans intérêt analytique. Le matériau historique suggère plutôt que chaque configuration d'émergence est singulière, et que ce qui manque à Babbage est identifiable au cas par cas.}. La leçon est plus générale~: les non-émergences ne sont pas des accidents extérieurs à la trajectoire technique~; elles rendent visible, en creux, la configuration d'éléments dont la rencontre permet à une capacité informationnelle de devenir infrastructure. Demande de coordination, support matériel stabilisé, organisation capable de porter l'usage, véhicule financier adapté, promesse socialement crédible~: cette configuration n'est pas une théorie causale (chacun de ces éléments dépend des autres) mais une grille heuristique que le matériau permet d'inférer.

Une seconde irrégularité interne au matériau force à préciser ce que la convergence morphologique cache. Calculatrice mécanique et tabulatrice mécanisent des opérations cognitives proches pour des organisations dont la fonction comptable, statistique ou administrative est comparable, et appartiennent à la même grande famille du calcul instrumenté. Pourtant l'une produit un marché fragmenté de plusieurs dizaines de fabricants, où Burroughs, Felt, Brunsviga, Comptometer, Monroe, Marchant et d'autres se partagent une clientèle qui possède ses machines, peut changer de fournisseur et opère hors d'un service vendeur~; l'autre produit une structure concentrée que la fusion CTR de 1911 puis l'IBM de Watson transforment en monopole effectif pendant un demi-siècle, fondé sur la location, le format propriétaire des cartes, la formation des opérateurs et le service de maintenance. L'écart interdit d'expliquer la structure de marché par la seule fonction technique. Il oblige à reconnaître que le mode juridique et économique d'appropriation du capital informationnel~--- vente ou location, possession ou abonnement, format ouvert ou propriétaire, machine séparée ou machine inséparable de ses consommables~--- peut être un déterminant de la structure de marché au même titre que les coûts fixes, les rendements croissants ou les externalités. Le mécanisme ne s'applique pas à toutes les configurations~: pour le câble sous-marin, l'ampleur du capital fixe et la spécificité des actifs ferment le choix, et l'on ne loue pas un câble transatlantique à des utilisateurs finaux. Mais lorsque la technique laisse plusieurs régimes d'usage techniquement disponibles, ce choix devient un déterminant économique indépendant des propriétés techniques, et c'est dans cette borne que le Ch.~\ref{ch:structures} aura à reprendre la dimension D3.

La seconde régularité concerne le travail. Dans chaque cas où une technique de l'information est introduite dans une organisation existante~--- la calculatrice mécanique dans les bureaux d'assurance (§\,\ref{sec:calculating}), la tabulatrice au \emph{Nautical Almanac} et au PLM (§\,\ref{sec:mechanography}), le commutateur automatique dans le réseau Bell (§\,\ref{sec:telephone})~---, le nombre de travailleurs se maintient ou augmente~; leur profil change~; le coût unitaire baisse par compression salariale plutôt que par suppression de postes. La mécanisation réorganise le travail au lieu de le remplacer~--- un résultat que \textcite{reinstaller_market_2001}, \textcite{davies_womans_1982}, \textcite{feigenbaum_answering_2024} et \textcite{daston_calculation_2018} documentent dans quatre registres différents (quantitatif, genre, emploi, organisation). Le seul cas qui produit une identification causale de ce mécanisme est celui de l'automatisation de la commutation téléphonique~: \textcite{feigenbaum_answering_2024}, exploitant la variation géographique du calendrier d'adoption dans près de trois mille villes américaines, montrent que les cohortes suivantes se réallouent vers des tâches nouvelles dans des industries qui n'employaient pas de secrétaires auparavant~--- du \og \emph{new work}\fg{} au sens de \textcite{autor_new_2024}. Mais ce mécanisme de réinstauration est conditionnel~: dans les villes frappées par la Grande Dépression, l'automatisation réduit l'emploi sans le compenser. La demande agrégée est une condition nécessaire pour que la restructuration l'emporte sur la suppression.

Cette régularité est le résultat le plus robuste du chapitre sur la dimension du travail. Mais elle ne dit pas tout. En relisant les cas traversés, trois observations apparaissent qui ne se réduisent pas à la question \og que fait la technique au travail~?\fg{} et qui suggèrent que le rapport entre les techniques de l'information et le travail est plus riche que ce que la formulation \og restructuration sans suppression\fg{} laisse entendre.

La première observation est que la technique informationnelle ne fonctionne qu'insérée dans un dispositif de travail intellectuel qu'elle ne peut pas remplacer. McCallum, sur l'\emph{Erie Railroad}, déploie six formes informationnelles dont le télégraphe n'est que la quatrième~; les cinq autres~--- le \emph{rule book}, la hiérarchie managériale, le reporting, la comptabilité analytique, la standardisation~--- sont du travail de conception codifié (§\,\ref{sec:telegraph}). Au \emph{Nautical Almanac}, Comrie estime que vingt à trente pour cent du temps du surintendant sont désormais consacrés à la conception des séquences dans lesquelles humains et machines doivent s'articuler~--- ce que \textcite{daston_calculation_2018} nomme l'\emph{analytical intelligence}. Bolle, au PLM, formule la même leçon dans une devise~: \og \emph{d'abord organiser, ensuite mécaniser}\fg{} (§\,\ref{sec:mechanography}). Dans ces trois cas, la machine est un composant d'un système dont le reste est du travail qualifié de conception~--- un travail que la mécanisation \emph{crée} au lieu de le réduire. Le chapitre~2 devra déterminer si ce coût de conception est un coût transitoire d'ajustement, comme le suggère \textcite{david_dynamo_1990} pour le parallèle dynamo/ordinateur, ou un coût permanent inhérent à toute mécanisation du traitement de l'information.

La seconde observation est que le stock de compétences existant ne subit pas passivement la technique~: il en conditionne la forme. L'arbitrage entre le système Morse (un fil, un opérateur formé au code) et le système Wheatstone (cinq fils, un opérateur non qualifié) n'est pas un choix purement technique~: c'est un arbitrage entre capital physique et capital humain, dans lequel la disponibilité de compétences détermine quel design s'impose (§\,\ref{sec:telegraph}). L'administration française du télégraphe fait construire par Bréguet un appareil électrique qui reproduit les mouvements du sémaphore optique pour ne pas avoir à requalifier les opérateurs~: la technologie adoptée est déformée par le stock de compétences en place. Et le métier Jacquard (§\,\ref{sec:era0}), en substituant du capital informationnel (les cartes perforées) au travail qualifié des tireurs de lacs, provoqua la résistance la plus célèbre de l'histoire des techniques~: les tisserands luddites (1811--1816) ne détruisaient pas des machines à vapeur mais des machines à information, des métiers dont les cartes perforées encodaient le savoir-faire qu'ils avaient passé des années à acquérir. Ce n'est pas la technique qui transforme le travail~: c'est le travail~--- sous forme de compétences disponibles ou indisponibles, défendues ou dévaluées~--- qui sélectionne, déforme ou rejette la technique. Le chapitre~2 devra instruire cette direction causale inverse avec les outils appropriés et déterminer si le stock de capital humain spécifique peut constituer une contrainte structurelle sur l'émergence d'un régime technologique~--- et pas seulement un facteur d'ajustement.

La troisième observation est que certaines tâches résistent à l'automatisation non par défaut technique mais parce qu'elles portent un savoir organisationnel implicite dont l'extraction exige la refonte du système entier. L'automatisation de la commutation téléphonique a pris soixante ans, de Strowger (1891) à la généralisation~; \textcite{feigenbaum_organizational_2024} ont montré que l'explication ne résidait ni dans la résistance syndicale ni dans un défaut de la technologie, mais dans le fait que la commutation interagissait avec la facturation, les réclamations, les renseignements et les urgences~--- une tâche \emph{intégrale} au sens de \textcite{ulrich_role_1995}, dont la modification impliquait la reconfiguration de l'ensemble du système Bell. Le chapitre~2 devra déterminer si ce type de résistance structurelle se retrouve dans les vagues d'automatisation ultérieures, et si la transgression de la frontière du savoir tacite par les modèles de langage~--- le \og paradoxe de Polanyi\fg{} de \textcite{autor_polanyi_2014}~--- modifie ou non cette donnée.

Le chapitre a enfin fait apparaître, sans le thématiser, un pattern biographique récurrent~: des individus formés dans un paradigme technique transfèrent leurs compétences vers le paradigme suivant par leur mobilité. Edison, opérateur de télégraphe, transpose la maîtrise des circuits électriques vers l'électrification (§\,\ref{sec:electrification}). Siemens, officier d'artillerie, transpose la connaissance des signaux électromagnétiques vers le télégraphe prussien (§\,\ref{sec:telegraph}). Shannon connecte l'algèbre de Boole au circuit de commutation (§\,\ref{sec:electrification}). La récurrence suggère que la diffusion des paradigmes techniques passe en partie par la mobilité des personnes qui en maîtrisent les composants~--- mais c'est une observation biographique, pas un résultat que le chapitre ait pu formaliser.

Ces quatre observations ne se substituent pas à la régularité centrale. Elles la complètent en montrant que le rapport entre les techniques de l'information et le travail est multidirectionnel~: la technique restructure le travail, mais elle crée aussi du travail de conception qu'elle ne remplace pas~; le travail, sous forme de compétences, sélectionne la technique~; et le savoir tacite porté par certaines tâches résiste à l'automatisation pendant des décennies. La littérature contemporaine sur l'automatisation raisonne presque exclusivement dans la première direction. Le chapitre~2, en mobilisant le cadre \emph{task-based} de \textcite{autor_skills_2003} et le modèle de \textcite{acemoglu_race_2018}, devra déterminer lesquelles de ces directions sont des phénomènes propres au dix-neuvième siècle et lesquelles sont des traits structurels du rapport ICT-travail qui persistent~--- ou se transforment~--- à l'ère de l'intelligence artificielle.

La troisième régularité concerne le rapport entre la technique de l'information et la demande qui l'appelle. La bidirectionnalité posée dans le préambule méthodologique~--- la technique transforme l'économie, mais l'économie appelle la technique~--- n'est pas symétrique à chaque période. Au dix-neuvième siècle, c'est surtout la demande qui tire~: la coordination militaire pour le télégraphe Chappe, la sécurité ferroviaire pour le \emph{block system}, le commerce pour le câble transatlantique, le recensement pour la tabulatrice. Au vingtième siècle, le basculement s'opère~: le modèle IBM (location et cartes exclusives) crée une dépendance au fournisseur~; Vail construit le réseau longue distance avant que la demande n'existe~; la Navy finance le wireless et le computing avant toute demande commerciale. Au vingt et unième siècle, \textcite{roussilhe_phase_2025} montre que l'offre ne se contente plus de précéder la demande~: elle doit la produire activement et continument, parce que la structure de coûts de l'industrie des semi-conducteurs exige que la puissance de calcul mise sur le marché chaque année soit \og consommée\fg{} pour maintenir les économies d'échelle. La flèche causale s'est inversée~--- et cette inversion a une conséquence énergétique directe que le chapitre a documentée mais ne peut pas modéliser.

L'anomalie est l'énergie. De §\,\ref{sec:era0} à §\,\ref{sec:convergence}, la dimension 0-TER apparaît dans chaque cas mais n'y constitue jamais un résultat mesurable. Le ratio entre la puissance du circuit informationnel et la puissance du système qu'il coordonne est de l'ordre de 1\,:\,2\,000 dans le cas ferroviaire. L'électricité représente 1,5\,\% du coût d'exploitation des tabulatrices au \emph{Nautical Almanac}. Le signal télégraphique tire son énergie de piles dont le coût n'apparaît nulle part dans les comptes de Western Union. La matérialité extractive~--- la gutta-percha, le cuivre, les poteaux~--- est réelle mais n'est documentée que rétrospectivement \parencite{tully_victorian_2009, fitzmaurice_material_2023}. L'énergie traverse le chapitre comme un fait sans observateur.


\subsection{Deux postulats qui cessent de tenir}\label{subsec:postulats}

Cette anomalie n'est pas un accident d'observation. Elle est produite par deux postulats implicites que le matériau historique n'a jamais contestés~--- jusqu'au cas contemporain.

Le premier postulat est que le coût marginal de l'information tend vers zéro. Le télégramme coûte à envoyer, mais l'information qu'il porte~--- le cours du coton à Liverpool~--- se diffuse gratuitement une fois connue. La calculatrice mécanique a un coût fixe d'achat, mais chaque addition supplémentaire ne coûte rien. Le logiciel Windows a un coût de développement, mais chaque PC supplémentaire qui l'exécute ne coûte rien à Microsoft en puissance de calcul~: il n'y a pas de coût marginal pour l'éditeur de logiciel \parencite{roussilhe_computing_2025}. Ce postulat tient de §\,\ref{sec:telegraph} à §\,\ref{sec:convergence}. Il cesse de tenir avec l'IA générative~: chaque inférence a un coût marginal positif, directement proportionnel à la consommation électrique par token. Pour la première fois dans l'histoire des techniques de l'information traversée par ce chapitre, l'énergie est un coût marginal de production du service informationnel.

Le second postulat est que l'efficacité par unité fonctionnelle s'améliore continument. Le coût par opération arithmétique baisse d'environ trois pour cent par an de 1850 à 1945 \parencite{nordhaus_two_2007}. L'efficacité énergétique par calcul double tous les dix-huit mois entre 1946 et 2000 \parencite{koomey_implications_2011}. Le \emph{scaling} de Dennard (§\,\ref{sec:electrification}) maintient la densité de puissance constante pendant que la miniaturisation progresse. Ce postulat cesse de tenir sur deux fronts~: le \emph{scaling} de Dennard atteint sa limite physique vers~2005, et l'empreinte environnementale par centimètre carré de fabrication de circuits intégrés, qui baissait depuis les années~1990, remonte à partir du n\oe{}ud de 7~nm \parencite{pirson_environmental_2022}. L'amélioration relative a cessé~; les impacts absolus croissent.

Ce qui est remarquable n'est pas que ces deux postulats cassent~--- toute trajectoire technique finit par rencontrer ses limites~--- mais qu'ils cassent simultanément. Tant que le coût marginal était nul, la question de l'efficacité énergétique ne se posait pas à l'utilisateur final~: Microsoft pouvait \og gaspiller les transistors\fg{} \parencite{gilder_coming_1995} sans que cela ne lui coûte quoi que ce soit. Tant que l'efficacité par composant s'améliorait, la consommation agrégée pouvait croître sans que le couplage ne devienne visible~: le \emph{scaling} de Dennard absorbait la croissance. Quand les deux cessent de tenir en même temps, le couplage entre information et énergie, invisible pendant deux siècles, apparaît. Le fait que personne n'ait croisé les séries de \textcite{nordhaus_two_2007} (coût par calcul) et de \textcite{koomey_implications_2011} (efficacité énergétique par calcul) n'est pas un oubli~; c'est le signe que les catégories d'analyse en vigueur ne contenaient pas la case.


\subsection{L'occultation comme pattern structurel}\label{subsec:occultation}

\textcite{ensmenger_environmental_2018} est le premier à constituer cette absence en objet d'étude. Dans son programme pour une histoire environnementale du \emph{computing}, il identifie trois axes~: les liens entre production de puissance de calcul et processus traditionnels d'extraction et de pollution~; la production de discours d'invisibilisation (la \og dématérialisation\fg{}, le \og cloud\fg{})~; l'asymétrie géographique entre zones d'extraction et zones de consommation. Mais ce qui fait la force de sa proposition, c'est que le deuxième axe~--- le discours~--- n'est pas un bruit à traverser pour atteindre les \og vrais faits\fg{} du premier axe. Il est constitutif~: c'est parce que le \emph{computing} est perçu comme immatériel que ses coûts matériels ne sont pas comptés, et c'est parce qu'ils ne sont pas comptés qu'ils peuvent croître sans rencontrer de résistance.

Le chapitre qui s'achève a documenté ce pattern bien au-delà du \emph{computing}. \textcite{belloc_telegraphie_1894} constate que la télégraphie est \og surtout remarquable par ses conséquences économiques\fg{} sans mentionner ses conséquences matérielles. D'Alembert juge le calcul \og plus laborieux que profond\fg{}, effaçant le travail du calcul comme objet digne d'attention. Gray ne voit pas la valeur commerciale du téléphone parce que ni la métrique du paradigme télégraphique~--- messages par fil, coût par mot~---, ni le modèle d'affaires de son employeur (Western Union, monopole télégraphique en place), ne contiennent la case (§\,\ref{sec:telephone}). \textcite{nordhaus_two_2007} mesure le coût du calcul sur deux siècles sans inclure l'électricité. Ce ne sont pas des erreurs~; ce sont des effets de cadrage. Chaque époque produit une grille qui capture certaines dimensions~--- le monopole naturel pour les régulateurs du télégraphe, la productivité du travail pour les historiens du calcul, l'intégration des marchés pour les économistes du commerce~--- et qui, structurellement, ne contient pas la case énergie-matière. L'occultation n'est pas produite par l'ignorance mais par les catégories d'analyse elles-mêmes.

\textcite{roussilhe_phase_2025} montre que ce pattern se renforce dans la période contemporaine. Les \og lois\fg{} qui organisent l'industrie des semi-conducteurs~--- la loi de Moore, les \emph{scaling laws} d'OpenAI, la \og loi de Huang\fg{}~--- ne sont pas des descriptions physiques~; ce sont des dispositifs de coordination sectorielle qui alignent les anticipations des fabricants, des éditeurs de logiciel et des investisseurs sur une feuille de route commune. Leur fonction est performative~: tant que l'industrie y croit, elle investit les sommes nécessaires pour les réaliser. Comme le formule Roussilhe, la loi de Moore est un \og baromètre de la croyance de l'industrie en elle-même\fg{}. La régulation du spectre électromagnétique (§\,\ref{sec:wireless}) procédait de la même structure~: \textcite{hazlett_rationality_1990} a montré que le Radio Act de~1927 n'était pas une correction de défaillance de marché mais un choix de distribution de rentes, organisé comme une réponse technique à un problème qui était en réalité politique. Les \og lois\fg{} technologiques et les régulations qui les accompagnent ne décrivent pas le monde~; elles le produisent~--- et ce qu'elles produisent inclut systématiquement l'invisibilité de la dimension énergétique.


\subsection{Ce que le chapitre a révélé et ce qu'il ne peut pas faire}\label{subsec:limites}

Le tableau~\ref{tab:ch1-questions} rassemble, pour chacune des dix dimensions et des deux forces transversales, ce que la traversée historique a permis d'établir et les questions qu'elle ouvre pour le chapitre~2. Trois niveaux sont distingués~: \textbf{★}~fait documenté indépendamment, \textbf{+}~argument explicite mais fondation incomplète, \textbf{?}~piste ouverte.

\begin{table}[p]
\caption{Les dix dimensions après la traversée historique~: acquis et questions ouvertes pour le chapitre~2.}\label{tab:ch1-questions}
\centering\scriptsize
\renewcommand{\arraystretch}{1.18}
\begin{tabularx}{\textwidth}{@{}cp{1.8cm}Xp{5.2cm}@{}}
\toprule
 & \textbf{Dimension} & \textbf{Ce que le ch.~1 établit} & \textbf{Questions pour le ch.~2} \\
\midrule
1 & Propriétés du bien & Le circuit télégraphique est rival~; l'information transportée ne l'est pas. La codifiabilité filtre les effets du commerce (Juhász~★). & ★ La codifiabilité module-t-elle \emph{toutes} les dimensions~? \newline ? Si l'IA rend codifiable presque tout, est-elle une ICT qui élimine le filtre~? \\[4pt]
2 & Rendements et concentration & Sept trajectoires institutionnelles convergent vers le monopole, par cinq mécanismes différents (★). Le résultat est stable~; le mécanisme varie. & ★ Qu'est-ce qui détermine le choix du mécanisme~? \newline + Trait des réseaux de communication, ou des infrastructures lourdes en général~? \\[4pt]
3 & Nature du capital & Le gradient isolé/intégré détermine irréversibilité et structure de marché (★ calculatrice vs tabulatrice). La forme du capital, pas la technologie, discrimine. & ★ Le mode d'appropriation (vente vs location+consommable) est-il le déterminant endogène de la structure de marché~? \newline + L'irréversibilité est-elle financière (Dixit-Pindyck) ou systémique (Hughes-Unruh)~? \\[4pt]
4 & Travail et nature du travail & La mécanisation réorganise le travail au lieu de le remplacer~--- quatre cas sur 150~ans (★). Conditionnalité~: la demande agrégée est nécessaire pour que la compensation opère (★ Feigenbaum-Gross). & ★ La conditionnalité à la demande agrégée tient-elle quand l'automatisation touche les tâches cognitives non-routinières~? \newline + Le travail de conception (\emph{analytical intelligence}) est-il un coût transitoire ou permanent~? \newline + Le stock de compétences conditionne-t-il la forme de la technologie adoptée~? \newline + Le savoir tacite des tâches intégrales constitue-t-il une résistance mesurable~? \\[4pt]
5 & Demande et préférences & La demande de calcul est dérivée et inélastique au prix de la machine (★). Le modèle «~service gratuit financé par la publicité~» naît avec la radio, pas avec Internet. & ? Effet de rebond par les tarifs~? \newline + La dissociation prix perçu / coût de production est-elle un trait constitutif de l'économie numérique~? \\[4pt]
6 & Signal de prix & L'information synchrone \emph{crée} le marché~--- Steinwender ★, quatre réplications. Effet conditionnel~: sans transport physique, pas d'arbitrage (Andrabi ★). & ★ La codifiabilité est-elle le seul filtre, ou la fonction ICT (TRANSMIT vs COMPUTE) en est-elle un autre~? \newline + Le seuil de McCallum (l'information comme condition de fonctionnement, pas simple facilitateur) se généralise-t-il~? \\[4pt]
7 & Coûts de transaction & La hiérarchie informationnelle émerge quand le prix ne coordonne pas (★ McCallum, six formes). La commutation est une tâche intégrale~: 60~ans pour automatiser (★ Feigenbaum). & + La frontière firme-marché est-elle déterminée par les techniques d'information disponibles~? \newline + La devise de Bolle (\og d'abord organiser, ensuite mécaniser~\fg{}) a-t-elle un équivalent contemporain~? \\[4pt]
8 & Financement & Pattern investissement $\rightarrow$ faillite $\rightarrow$ recapitalisation documenté pour le câble (★), la railway mania, le dot-com. L'État crée le marché avant le marché (★ Census, Navy, Convention). & + Le cycle Perez est-il propre aux réseaux de communication ou commun aux infrastructures lourdes~? \newline ? Le financement par la promesse stratégique (pas la rentabilité) est-il propre aux technologies de souveraineté~? \\[4pt]
9 & Géographie & La géographie du réseau reflète les rapports de puissance (★ câbles coloniaux, concessions). Domination britannique sur les matériaux du câble (★). & + La contrainte matérielle (cuivre, gutta-percha) préfigure-t-elle les terres rares et le lithium~? \newline ? La localisation des DC suit-elle l'énergie, le pouvoir, ou les rendements croissants~? \\[4pt]
10 & Innovation & Le système métrique du paradigme en place rend invisible la valeur des innovations d'un autre registre~--- Gray, Marconi (★). Verrouillage de standard sur des millénaires (sexagésimal, VN, carte 80~col.). & ★ L'occultation par le système métrique se reproduit-elle à chaque vague~? \newline + Le facteur limitant de l'innovation migre-t-il du talent vers l'énergie dans la période contemporaine~? \\[4pt]
\midrule
0-TER & Énergie-matière & L'énergie traverse le chapitre comme un fait sans observateur. Ratio 1:2\,000 (circuit informationnel / système coordonné). Deux postulats (coût marginal nul, efficacité croissante) cassent simultanément. & ★ Le ratio évolue-t-il~? À quel seuil le couplage devient-il visible~? \newline + L'occultation est-elle produite par les catégories d'analyse ou par les ordres de grandeur~? \\[4pt]
0-GOV & Souveraineté & La forme du réseau reflète l'intérêt de l'acteur qui le finance~--- de Chappe aux PPAs nucléaires (★). & + La souveraineté numérique prolonge-t-elle le même pattern~? \\
\bottomrule
\end{tabularx}
\renewcommand{\arraystretch}{1.0}
\end{table}

\begin{table}[p]
\caption{Matrice dimensions $\times$ périodes~: trajectoire de l'occultation énergétique.}\label{tab:synoptique-ch1}
\centering\scriptsize
\renewcommand{\arraystretch}{1.12}
\setlength{\tabcolsep}{3pt}
\begin{tabular}{@{}llcccccccccccc@{}}
\toprule
 & \textbf{Période} & \rotatebox{70}{\textbf{D1}} & \rotatebox{70}{\textbf{D2}} & \rotatebox{70}{\textbf{D3}} & \rotatebox{70}{\textbf{D4}} & \rotatebox{70}{\textbf{D5}} & \rotatebox{70}{\textbf{D6}} & \rotatebox{70}{\textbf{D7}} & \rotatebox{70}{\textbf{D8}} & \rotatebox{70}{\textbf{D9}} & \rotatebox{70}{\textbf{D10}} & \rotatebox{70}{\textbf{0-TER}} & \rotatebox{70}{\textbf{0-GOV}} \\
\midrule
§\ref{sec:era0} & Avant 1794 & $\bullet$ & --- & $\circ$ & $\bullet$ & $\circ$ & --- & $\circ$ & --- & $\circ$ & $\bullet$ & --- & $\circ$ \\
§\ref{sec:telegraph} & 1794--1880 & $\bullet$ & $\bullet$ & $\circ$ & $\bullet$ & $\bullet$ & $\bullet$ & $\circ$ & $\bullet$ & $\bullet$ & $\bullet$ & --- & $\bullet$ \\
§\ref{sec:calculating} & 1820--1914 & $\circ$ & $\circ$ & $\bullet$ & $\bullet$ & $\circ$ & --- & --- & $\circ$ & --- & $\circ$ & --- & --- \\
§\ref{sec:mechanography} & 1890--1940 & $\circ$ & $\bullet$ & $\bullet$ & $\bullet$ & $\circ$ & --- & $\circ$ & $\bullet$ & --- & $\bullet$ & --- & $\circ$ \\
§\ref{sec:telephone} & 1876--1920 & $\circ$ & $\bullet$ & $\circ$ & $\bullet$ & $\bullet$ & $\circ$ & $\bullet$ & $\circ$ & $\circ$ & $\bullet$ & --- & $\bullet$ \\
§\ref{sec:wireless} & 1895--1930 & $\circ$ & $\bullet$ & --- & $\circ$ & $\bullet$ & $\bullet$ & --- & $\circ$ & --- & $\bullet$ & --- & $\bullet$ \\
§\ref{sec:electrification} & 1880--1930 & --- & $\bullet$ & $\bullet$ & $\circ$ & $\bullet$ & --- & $\bullet$ & $\bullet$ & --- & $\bullet$ & $\nearrow$ & $\bullet$ \\
§\ref{sec:computing} & 1945--1975 & --- & $\bullet$ & $\bullet$ & $\circ$ & $\bullet$ & --- & --- & $\bullet$ & --- & $\bullet$ & $\circ$ & $\bullet$ \\
§\ref{sec:convergence} & 1975--2007 & $\circ$ & $\bullet$ & $\bullet$ & $\circ$ & $\bullet$ & $\bullet$ & $\bullet$ & $\bullet$ & $\circ$ & $\bullet$ & $\bullet$ & $\circ$ \\
§\ref{sec:contemporain} & 2007-- & $\circ$ & $\bullet$ & $\bullet$ & $\nearrow$ & $\bullet$ & $\circ$ & $\circ$ & $\bullet$ & $\bullet$ & $\bullet$ & $\bullet$ & $\bullet$ \\
\midrule
\multicolumn{2}{@{}l}{\textbf{Synthèse}} & \multicolumn{10}{c}{\emph{Dimensions stables~: D2, D5, D8, D10}} & $\boldsymbol{\nearrow\nearrow}$ & \\
\bottomrule
\end{tabular}
\renewcommand{\arraystretch}{1.0}
\par\medskip
\footnotesize
\textbf{Légende}~: $\bullet$~=~dimension centrale dans cette période~; $\circ$~=~dimension présente mais secondaire~; ---~=~dimension non pertinente ou non documentée~; $\nearrow$~=~dimension en émergence.
\par\smallskip
\textbf{Lecture de la colonne 0-TER}~: La dimension énergie-matière passe de «~---~» (non pertinente) dans les sections §\ref{sec:era0}--§\ref{sec:wireless}, à «~$\nearrow$~» (émergente) en §\ref{sec:electrification} avec l'électrification du substrat, à «~$\circ$~» (présente mais secondaire) en §\ref{sec:computing} où l'énergie est une contrainte d'ingénierie mais pas un coût économique visible, à «~$\bullet$~» (centrale) à partir de §\ref{sec:convergence}. Cette trajectoire illustre l'\emph{occultation progressive puis la réémergence} de la contrainte énergétique dans l'analyse des techniques de l'information.
\end{table}

Le chapitre a révélé une structure. Dix dimensions économiques (D1--D10) et deux forces transversales (0-TER, 0-GOV) ont émergé de la confrontation séquentielle des cas. Des hypothèses abductives ont été formulées~: le mode d'appropriation du capital, et non la technologie, détermine la structure de marché~; la mécanisation réorganise le travail au lieu de le remplacer~; la technique ne se déploie pas sans un véhicule de financement adapté à la structure des coûts et des bénéfices. Un gradient d'occultation de la dimension énergie-matière a été identifié, produit non par l'ignorance mais par les catégories d'analyse en vigueur à chaque époque. Et deux postulats implicites~--- le coût marginal nul de l'information, l'amélioration continue de l'efficacité~--- ont été nommés parce que le cas contemporain les a fait voler en éclats.

\dimmark{D10}{*}%
Un parallèle disciplinaire récent mérite d'être nommé ici par souci de positionnement bibliographique. \textcite{johnstone_deep_2026}, dans la revue \emph{Technological Forecasting and Social Change}, ont mobilisé le cadre des \emph{Deep Transitions} développé par \textcite{schot_deep_2018} pour analyser l'évolution du «~méta-régime numérique~» sur la même période 1937--2024 que celle parcourue par le présent chapitre, et la séquence de phases qu'ils en dérivent~--- gestation, irruption, frenzy, turning point, synergy, maturity~--- recouvre une partie des mêmes faits empiriques. Les deux approches ne se substituent pas l'une à l'autre. Leur cadre est évolutionniste et historique au sens de la sociologie des transitions sociotechniques (STS), construit autour des notions de règles dominantes, de règles précurseurs et de méta-régime~; le présent chapitre est dans le registre de l'induction analytique au sens de \textcite{znaniecki_method_1934}, de la causalité par cas nécessaire ou suffisante au sens de \textcite{trannoy_causalite_2025}, et de l'occultation matérielle au sens de \textcite{ensmenger_environmental_2018}. Le chapitre mobilise ponctuellement le récit de Johnstone et collègues là où il éclaire des points précis de l'observation par cas (§\ref{sec:electrification}, §\ref{sec:computing}, §\ref{sec:convergence}, §\ref{sec:contemporain}), sans transposer leur cadre théorique~: l'objet du présent chapitre est de produire un matériau brut sur lequel le cadre disciplinaire de la suite de la thèse (économie de la croissance et des techniques génériques, comptabilité énergétique, modélisation IMACLIM) puisse opérer, non d'opérationnaliser une approche STS-évolutionniste sur le même matériau. La distinction des registres importe aussi pour éviter une redondance interprétative~: ce que \textcite{johnstone_deep_2026} décrivent comme \emph{turning point} (2001--2008) apparaît dans le présent chapitre comme une rupture de couplage énergétique permanent (§\ref{subsec:turning_point})~; ce qu'ils nomment \emph{frenzy} (1996--2000) apparaît comme une résolution différée du modèle économique de la commercialisation d'Internet (§\ref{subsec:dotcom})~; ce qu'ils caractérisent comme phase de \emph{synergy} contemporaine (2008--) est décomposé dans le présent chapitre en cloud, store, GenAI, blockchain, IoT et régulation fragmentée (§\ref{sec:contemporain}). Les éclairages se complètent sans se recouvrir.

Mais le chapitre ne peut pas aller plus loin. L'induction identifie des régularités~; elle ne les explique pas. L'abduction génère des hypothèses~; elle ne les teste pas. Le constat que sept pays convergent vers le monopole télégraphique n'est pas un mécanisme~; le constat que l'occultation de 0-TER se reproduit à chaque cas n'est pas une cause. Pour formaliser les mécanismes qui produisent ces régularités~--- et pour déterminer lesquels sont propres aux techniques de l'information et lesquels sont communs à toutes les infrastructures lourdes~---, il faut changer d'instrument. L'analyse économique systématique du chapitre~2 prend les dix dimensions que ce chapitre a fait émerger et les déploie comme des canaux causaux, chacun avec ses mécanismes, ses conditions de validité et ses équations. L'analyse énergie-matière du chapitre~3 prend la dimension transversale 0-TER et la quantifie fonction par fonction~--- transmission, stockage, calcul~--- pour mesurer ce que l'occultation cachait. Le chapitre qui s'achève a rendu les deux suivants nécessaires~; il ne pouvait pas les remplacer.